\section{Analytical Backbone: What Is Proved, What Is Borrowed, What Is Computed}
\label{sec:analysis}

This section fixes the mathematics the experiments rest on. Each item is labelled by status: \emph{foundation} (standard, restated for self-containment, no priority claimed), \emph{borrowed tool} (prior art, attributed), \emph{proved here}, or \emph{computed} --- and, for computed quantities, whether the number is an interval enclosure or a high-accuracy floating-point evaluation. Three error objects recur and are never interchanged: the memory-kernel error, the induced response-operator error, and the fixed-input discrepancy (Section~\ref{subsec:semantics}).

\subsection{Result I: exact structural separation (foundation)}
The collocated Caputo transfer of the linearisation at $(x^*,y^*(A))$ carries an $O(\omega^{-\alpha})$ high-frequency signature. No finite rational latent transfer (integer-order rolloff) and no strictly proper retarded-delay model ($O(\omega^{-1})$) reproduces it exactly on an open frequency set. The mechanism --- a branch point of $s^\alpha$ at the origin against the meromorphy of the rival classes --- is classical fractional-systems material~\cite{F02,F03,F08}; no novelty is claimed for it. It is stated because it delimits what the later results can and cannot mean.

\begin{theorem}[Structural separation]
\label{thm:T4}
Let $CB\neq0$ and $\alpha\in(0,1)$ with $\alpha\notin\mathbb{Q}$ in the sense that no effectively present exponent of $s^\alpha$ is an integer multiple of $1/\alpha$. Then $G_\alpha$ is not meromorphic in any neighbourhood of $s=0$, hence differs from every finite rational transfer and from every strictly proper retarded-delay transfer.
\end{theorem}

The statement concerns exact equality of transfers on open frequency sets. It promises nothing about experiments with finite bandwidth, finite horizon $T<\infty$ and noise $\sigma>0$. That regime is governed by Lemma~\ref{lem:soe} and its corollary.

\subsection{Result II: positive exponential-sum approximation (borrowed tool)}
\label{subsec:soe}
The Caputo kernel $k_\alpha(t)=t^{\alpha-1}/\Gamma(\alpha)$ has the Stieltjes representation $k_\alpha(t)=c_\alpha\int_0^\infty\lambda^{-\alpha}e^{-\lambda t}\,d\lambda$ with $c_\alpha=\sin(\pi\alpha)/\pi$. Discretising it by the trapezoidal rule on a logarithmic grid $\lambda_k=e^{kh}$ and truncating produces an exponential sum with \emph{positive} weights and rates.

\begin{quote}
\textbf{Attribution.} This construction is prior art. Positive exponential-sum approximation of $t^{-\beta}$ by quadrature of an integral representation, including the positivity of the weights and the root-exponential order in the number of terms, is due to McLean~\cite{K07}, building on the exponential-sum approximation theory of Beylkin and Monz\'on~\cite{K03,K04}; the same construction underlies the fast Caputo evaluation of Jiang et al.~\cite{K05}. The quadrature estimate we invoke is the analytic-strip trapezoidal bound~\cite{K08}, and root-exponential order is the expected optimal order for an algebraic branch point~\cite{K09}. We claim no new approximation method, no new rate law, and no priority for positive weights.
\end{quote}

What we add is narrow and technical: the published estimates are stated on a compact interval $[\delta,T]$ with $\delta>0$, excluding the singularity at the origin, whereas the discrimination argument of Section~\ref{subsec:testing} needs the convolution norm $L^1(0,T)$ \emph{including} $t=0$. The strip bound is $t$-integrable there because $\alpha>0$, so no excision is required.

\begin{lemma}[Endpoint-inclusive $L^1$ estimate and exact horizon scaling]
\label{lem:soe}
Let $E_m^+(\alpha,T)$ denote the least $L^1(0,T)$ error of an $m$-term exponential sum with positive weights and rates. Then
\begin{equation}
E_m^+(\alpha,T)=T^\alpha E_m^+(\alpha,1),
\label{eq:scaling}
\end{equation}
and for every $c<\pi\sqrt{\alpha(1-\alpha)}$ there exists $A(\alpha,c)<\infty$ with
\begin{equation}
E_m^+(\alpha,T)\le A(\alpha,c)\,T^\alpha e^{-c\sqrt m}\qquad(m\ge1).
\label{eq:rootexp}
\end{equation}
Consequently $m(\varepsilon)=O(\log^2(1/\varepsilon))$ terms suffice for accuracy $\varepsilon$.
\end{lemma}

\begin{proof}
Equation~\eqref{eq:scaling} is exact: substituting $t=Ts$ and using $k_\alpha(Ts)=T^{\alpha-1}k_\alpha(s)$ turns the $L^1(0,T)$ error into $T^\alpha$ times the $L^1(0,1)$ error, and $(c_j,\lambda_j)\mapsto(T^{1-\alpha}c_j,T\lambda_j)$ is a bijection of the admissible positive mixtures. For~\eqref{eq:rootexp} put $\lambda=e^x$, so $k_\alpha(t)=c_\alpha\int_{\mathbb R}g(x,t)\,dx$ with $g(x,t)=e^{(1-\alpha)x}e^{-te^x}$. For $|\operatorname{Im}x|\le d<\pi/2$ one has $|g|\le e^{(1-\alpha)x}e^{-te^x\cos d}$, whence the strip norm obeys $N(g(\cdot,t),d)\le2(\cos d)^{\alpha-1}\Gamma(1-\alpha)t^{\alpha-1}$. Integrating the trapezoidal estimate of~\cite{K08} over $t\in(0,1)$ and using $\int_0^1t^{\alpha-1}dt=1/\alpha$ together with $c_\alpha\Gamma(1-\alpha)=1/\Gamma(\alpha)$ bounds the discretisation error by $2\Gamma(\alpha+1)^{-1}(\cos d)^{\alpha-1}e^{-2\pi d/h}(1-e^{-2\pi d/h})^{-1}$; the truncation tails contribute $O(e^{-\alpha M_+h})$ and $O(e^{-(1-\alpha)M_-h})$. Equalising the three exponents gives $m\approx2\pi d\,h^{-2}[\alpha(1-\alpha)]^{-1}$ and exponent $\sqrt{2\pi d\,\alpha(1-\alpha)m}$, which tends to $\pi\sqrt{\alpha(1-\alpha)m}$ as $d\uparrow\pi/2$ while $(\cos d)^{\alpha-1}\to\infty$. Hence any $c$ strictly below the limit is attained with a finite constant.
\end{proof}

The quantifier order in~\eqref{eq:rootexp} is not cosmetic: the constant diverges as $c$ approaches $\pi\sqrt{\alpha(1-\alpha)}$, so the boundary rate is \emph{not} established, and whether it is attained or optimal is open. Equation~\eqref{eq:scaling} is verified to machine precision over a 285-cell sweep ($\alpha\in\{0.30,\dots,0.95\}$, $T\in\{1,10,100\}$, $m\le128$): the measured ratio $E(T)/E(1)$ equals $T^\alpha$ to all printed digits, and the relative error is $T$-invariant. Over the same sweep, \eqref{eq:rootexp} holds at $c=\pi\sqrt{\alpha(1-\alpha)}$ with constants between $0.29$ and $11.1$, and is attained to within a factor $1.00$ at $\alpha=0.70$, $m=128$.

\begin{corollary}[Finite-horizon closure of the latent hierarchy]
\label{cor:closure}
For every $T<\infty$ and $\varepsilon>0$ some finite positive exponential mixture lies within $\varepsilon$ of $k_\alpha$ in $L^1(0,T)$. Hence, absent a declared bound on the latent dimension, no fixed finite-horizon experiment can guarantee a uniform positive separation from the whole latent hierarchy: discrimination is meaningful only relative to a budget $m$.
\end{corollary}

Results I and II together fix the geometry: exact separation (Theorem~\ref{thm:T4}) coexists with $\varepsilon$-approximation at finite $T$ (Corollary~\ref{cor:closure}). In the benchmark of Section~\ref{sec:benchmark} the latent class is capped at $m=3$; in the response-level computation of Section~\ref{subsec:state} the budget runs to $m=128$.

\subsection{Semantics: three error objects that must not be conflated}
\label{subsec:semantics}
Write $G_\alpha(s)=C(s^\alpha I-J)^{-1}B$ for the exact prey-to-prey transfer and $G_m$ for a strictly proper finite-state surrogate. Three distinct quantities appear below.
\begin{itemize}
\item[(O1)] \emph{Memory-kernel error} $E_m=\|k_\alpha-k_m\|_{L^1(0,T)}$: a statement about the fractional operator alone, independent of $J$, $B$, $C$. This is the object of Lemma~\ref{lem:soe}.
\item[(O2)] \emph{Induced response-operator error} $\widehat E_m^{\rm state}=\inf_{\mathcal G_m}\|g_\alpha-g_m\|_{L^1(0,T)}$, the $L^1$ distance between impulse responses, i.e.\ the $L^\infty\!\to\!L^\infty$ operator norm of the difference of input--output maps. It is a \emph{minimax} quantity: an infimum over rivals of a supremum over inputs.
\item[(O3)] \emph{Fixed-input discrepancy} $\|(g_\alpha-g_m)*u\|/\|u\|$ for one specific $u$.
\end{itemize}
By construction $\text{(O3)}\le\text{(O2)}$ for every admissible $u$. There is no inequality in either direction between (O1) and (O2), and no constant relating them: computed from the same surrogate family, the plant \emph{attenuates} the kernel error monotonically, by a factor $1.009$ at $m\approx5$ falling to $0.675$ at $m=128$. Kernel-level numbers are therefore never propagated into response-level claims anywhere in this paper.

Two consequences are used later. First, because (O2) is a supremum over inputs, the testing bound of Section~\ref{subsec:testing} is valid uniformly over the admissible input class, and conservative. Second, and for the same reason, $\widehat E_m^{\rm state}$ is an \emph{obstruction certificate and not a design objective}: the frequency band that maximises $|G_\alpha-G_m|$ (here $\omega\approx0.4$) selects the multiscale waveform, which excites $87\%$ of the worst case at $m=32$ and is nonetheless the \emph{least} informative design in the four-class benchmark (Table~\ref{tab:safe-search}). Maximising one pairwise discrepancy does not maximise multi-class accuracy.

\subsection{Result III: testing obstruction at the response-operator level (proved here)}
\label{subsec:testing}
Fix the observation model: the prey channel is sampled at $n$ times in $(0,T]$ and corrupted by additive Gaussian noise with common covariance $\sigma^2I$ under both hypotheses. The two hypotheses are ``the response is $g_\alpha$'' and ``the response is $g_m$'', with means $\mu_\alpha$, $\mu_m$ obtained by convolving the respective response kernels with the same admissible input $u$. Because the covariance is common and known, the minimax error of the two-point test is available in closed form and no inequality is needed.

\begin{theorem}[Exact two-point obstruction under a response-operator budget]
\label{thm:T20}
Let $\|u\|_2\le B$ and let $d=\|\mu_\alpha-\mu_m\|_2/\sigma$. Then the minimax error of any test between the two hypotheses satisfies
\begin{equation}
P_e^*=\Phi\!\left(-\tfrac{d}{2}\right),\qquad
d\le\frac{\widehat E_m^{\rm state}B}{\sigma},
\label{eq:exact2pt}
\end{equation}
by Young's inequality $\|(g_\alpha-g_m)*u\|_2\le\|g_\alpha-g_m\|_{L^1}\|u\|_2$. Hence
\begin{equation}
P_e^*\ \ge\ \Phi\!\left(-\frac{\widehat E_m^{\rm state}B}{2\sigma}\right)\ \longrightarrow\ \tfrac12
\quad\text{as }\widehat E_m^{\rm state}\to0\text{ with }B,\sigma\text{ fixed.}
\label{eq:floor}
\end{equation}
\end{theorem}

Three points of hygiene. The bound is routed through the $L^1$ norm of the response kernel, never through $\|k_\alpha-k_m\|_{L^2}$, which is infinite for $\alpha\le1/2$ and would make any such statement vacuous on part of the declared range. The exact bound~\eqref{eq:exact2pt} is the primary statement; the Pinsker bound $\tfrac12(1-\sqrt{{\rm KL}/2})$ is reported alongside in Table~\ref{tab:state-approx} only for comparability with the literature, and is uniformly looser --- by $0.034$ at $m=4$.

\subsection{Result IV: finite-state prey-response approximation (computed)}
\label{subsec:state}
Results I--III concern operators. Result IV addresses the measured object, the prey-to-prey response $g_\alpha$. The factorisation $(s^\alpha I-J)^{-1}=(sI-J)^{-1}s^{-\alpha}$ is false and is not used. Instead $g_\alpha$ is split by contour deformation into its principal-sheet pole pair --- realised \emph{exactly} by a stable two-state integer-order block --- plus a branch-cut remainder, smooth on $(0,\infty)$, which is the only part requiring approximation.

\begin{theorem}[Certified finite-state approximation of the prey response]
\label{thm:T23}
For the displayed Matignon-stable cells and $m\in\{4,8,16,32\}$ there exists an explicit stable finite-state surrogate $g_m$ with $\|g_\alpha-g_m\|_{L^1(0,T)}\le\widehat E_m^{\rm state}$, the bound being an outward-rounded interval enclosure over the admissible pole region and $L^1$ mass cap.
\end{theorem}

The values of $\widehat E_m^{\rm state}$ for $m\le32$ (Table~\ref{tab:state-approx}) are \textbf{interval enclosures}. Table~\ref{tab:state-approx-large} continues the computation to $m=64$ and $m=128$ using one explicit positive-SOE surrogate rather than the optimised class; those numbers are \textbf{high-accuracy floating-point evaluations with a two-mesh quadrature-error estimate}, not enclosures, and because that surrogate is one admissible element rather than the infimum they are upper bounds on $\widehat E_m^{\rm state}$ --- so the floors they induce through~\eqref{eq:floor} are valid and conservative. The distinction is maintained wherever these numbers are used.

Combining Theorem~\ref{thm:T20} with the computed values gives the operative conclusion: the testing floor rises from $0.374$ at $m=4$ to $0.4996$ at $m=64$ and $0.49999$ at $m=128$. At attainable latent orders the finite-state rival is indistinguishable from the fractional model under the declared protocol, for \emph{every} input in the budget. This is the obstruction that the waveform search of Section~\ref{subsec:safe-search} does not remove, and it is separate from the safety question entirely.
